% Non-technical homeowner advisory (~1.5 pages).
% Compile: pdflatex homeowner_advisory_letter.tex
% Style: clear public prose; short sentences; no equations.
\documentclass[11pt]{article}
\usepackage[margin=0.9in]{geometry}
\usepackage{setspace}
\pagestyle{empty}
\setstretch{1.12}
\begin{document}

\begin{center}
{\Large\bfseries Smart Climate Control: How Predictive Pre-Cooling\\
Lowers Your Electric Bill Without Sacrificing Comfort}\\[0.6em]
{\small A practical note for homeowners \quad Team GoHiMCM \quad HiMCM 2018 Problem B}
\end{center}

\vspace{0.8em}
\noindent
Dear Homeowner,

\medskip
\noindent
Afternoon power is often the most expensive power of the day.
Many homes still cool only when the room already feels hot.
That habit pushes your compressor into the high-price window.
It also makes the unit start and stop many times.

\medskip
\noindent
Our student team studied a simple idea.
Cool the house a little earlier.
Let the walls hold that cool air.
Then ease the cooling when rates jump.
We call this predictive pre-cooling.
In our locked summer-day test, the predictive controller cut the bill by about
\textbf{4\%} and cut peak-period cooling energy by about \textbf{16\%}.
Compressor starts fell from \textbf{42} to \textbf{3} on that day.

\medskip
\noindent
You do not need a research lab to use the same logic.
Use these three household rules.

\medskip
\noindent\textbf{Rule 1 --- Pre-cool two hours before the weekday peak.}
If peak rates begin near 14:00, start a deeper cool near 12:00.
Aim for a slightly cooler house before the expensive window.
The walls then act like a thermal battery.
You buy more cooling when power is cheaper.

\medskip
\noindent\textbf{Rule 2 --- Raise the setpoint about $1.5^\circ\mathrm{C}$ from 14:00 to 19:00.}
During the peak window, let the house drift warmer by a small amount.
A brief rise toward the mid-$20^\circ\mathrm{C}$ range is often tolerable.
Your walls release stored coolth.
The compressor works less when each kilowatt-hour costs more.
In our family negotiation test, elders kept a floor near $21.5^\circ\mathrm{C}$.
Parents accepted a short float near $25^\circ\mathrm{C}$ on peak.
Comfort stayed the priority.
The bill still fell.

\medskip
\noindent\textbf{Rule 3 --- Keep the outdoor unit clean so COP stays high.}
A dirty outdoor coil forces the system to draw more electricity.
Wash leaves and dust from the condenser on a regular schedule.
Replace clogged filters.
Efficient hardware multiplies every smart schedule.

\medskip
\noindent
These rules work with many modern thermostats.
If you own a learning thermostat, set a weekday pre-cool block.
If you use a simple programmable unit, enter the same schedule by hand.
If several people share the home, agree on a comfort band first.
Protect anyone who is cold-sensitive.
Then allow a small peak float for everyone else.

\medskip
\noindent
Smart climate control is not about living in a cold house all day.
It is about moving cooling into cheaper hours.
It is about fewer short cycles.
It is about keeping comfort while the tariff curve spikes.

\medskip
\noindent
Sincerely,\\
GoHiMCM AeroTherm Team

\vspace{1.0em}
{\small
Evidence for the numeric claims above is archived in locked local simulation
files (\texttt{bang\_vs\_mpc\_summary.json},
\texttt{agent\_thermostat\_dialogue.json}).
}

\end{document}
